\chapter{What a Trading Firm Does}\label{ch:m1:what-a-trading-firm-does}

At 10:31:04 one hundred shares of a large company change hands at \$50.01.
The buyer is a dentist in Ohio who tapped a button on a phone; the seller is
a firm in a data centre that has never heard of her and held the shares for
eleven seconds. Between them, a broker, an exchange and a clearing house each
handled the order, and each was paid. The print lasted a microsecond on the
tape. This chapter is about who these firms are, how each of them earns a
living from that print, and why none of them is doing the same job.

\section{Five business models}

Every firm in this book lives from trading, but ``trading firm'' covers at
least five different businesses. They differ in three things: whose money is
at risk, who pays them, and how long they hold a position.

\begin{definition}[Market maker]\label{def:m1:what-a-trading-firm-does:market-maker}
A \emph{market maker}\index{market maker} is a firm that continuously offers
to buy and to sell the same instrument, at a lower price to buy (its
\emph{bid}) than to sell (its \emph{ask}), and trades with whoever accepts.
It risks its own capital and aims to end each day with little or no
position.
\end{definition}

\begin{definition}[Bid--ask spread and mid price]\label{def:m1:what-a-trading-firm-does:spread}
For a best bid $b$ and a best ask $a > b$, the
\emph{bid--ask spread}\index{bid--ask spread} is $s = a - b$ and the
\emph{mid price}\index{mid price} is $m = \tfrac12(a+b)$. The
\emph{half-spread} $s/2$ is what one market order pays relative to the mid.
\end{definition}

\begin{example}[The dentist's order]\label{ex:m1:what-a-trading-firm-does:dentist}
The market is $49.99$ bid, $50.01$ ask: $s = \$0.02$, $m = \$50.00$, and the
spread is $0.02/50.00 = \qty{4}{\bp}$ of the price. The dentist's market
order buys at the ask and pays the half-spread, \$0.01 a share, that is
\$1.00 on her hundred shares. That dollar is the \omterm{def:m1:what-a-trading-firm-does:market-maker}{market maker}'s gross revenue
on the trade --- before the price has had a chance to move.
\end{example}

\begin{definition}[Proprietary trading firm]\label{def:m1:what-a-trading-firm-does:prop}
A \emph{proprietary trading firm}\index{proprietary trading firm} trades
only its owners' capital, has no clients and charges no fees: its entire
revenue is its trading profit. Most electronic \omterm{def:m1:what-a-trading-firm-does:market-maker}{market makers} are proprietary
firms; many proprietary firms also run strategies that take liquidity
instead of providing it.
\end{definition}

\begin{definition}[Hedge fund and multi-manager platform]\label{def:m1:what-a-trading-firm-does:hedge-fund}
A \emph{hedge fund}\index{hedge fund} is a privately offered investment
vehicle that trades outside investors' capital with few constraints on
instruments, leverage or short selling, and is paid a fee on the assets plus
a share of the profits. A \emph{multi-manager platform}\index{multi-manager platform}
is a hedge fund organised as many independent teams (``pods''), each given
capital and tight risk limits by a central allocator that nets and manages
the sum.
\end{definition}

\begin{definition}[Asset manager and assets under management]\label{def:m1:what-a-trading-firm-does:asset-manager}
An \emph{asset manager}\index{asset manager} invests clients' money under a
written mandate --- typically to track or to beat a benchmark --- and is paid
a percentage of the \emph{assets under management}\index{assets under management}
(AUM), the market value of the portfolios it runs.
\end{definition}

The fifth model is the bank's trading desk: a dealer that makes prices to
clients in exchange for their business, studied in \cref{ch:m1:the-sell-side}.
\Cref{ex:m1:what-a-trading-firm-does:table} puts the five side by side.

\begin{example}[Five firms, one table]\label{ex:m1:what-a-trading-firm-does:table}
\begin{center}\small
\begin{tabular}{@{}lllll@{}}
\hline
 & Whose capital & Paid by & Holding period & Main cost \\
\hline
\omterm{def:m1:what-a-trading-firm-does:market-maker}{Market maker} & owners' & the spread & seconds to hours & technology \\
Proprietary firm & owners' & trading profit & seconds to weeks & technology, people \\
\omterm{def:m1:what-a-trading-firm-does:hedge-fund}{Hedge fund} & investors' & fees on assets and profits & days to years & people \\
\omterm{def:m1:what-a-trading-firm-does:asset-manager}{Asset manager} & clients' & fee on assets & months to years & distribution \\
Bank desk & shareholders' & spread and client business & minutes to months & capital \\
\hline
\end{tabular}
\end{center}
The rest of the series keeps returning to this table: a technique that is
central to one row (speed, for the first) is irrelevant to another (the
fourth).
\end{example}

\begin{remark}[Where ``two and twenty'' comes from]\label{rem:m1:what-a-trading-firm-does:jones}
The \omterm{def:m1:what-a-trading-firm-does:hedge-fund}{hedge fund}'s fee structure is as old as the \omterm{def:m1:what-a-trading-firm-does:hedge-fund}{hedge fund}. Alfred Winslow
Jones, who started the first one in 1949 with \$100\,000, took 20\% of the
profits and at first nothing else; an annual fee on assets came later, and
``2 and 20'' --- 2\% of assets, 20\% of profits --- became the reference from
which every fund now negotiates downward or, for the most sought-after,
upward.
\end{remark}

\begin{dated}{2026-09}{How large the industry is}\label{dat:m1:what-a-trading-firm-does:aum}
Professionally managed assets worldwide were about \$147 trillion at the end
of 2025, up from \$128 trillion a year earlier, by the most widely quoted
industry count. Most of that year's growth came from rising markets, not from
new money.
\end{dated}

\section{Where the money comes from}

Strip away the vocabulary and a trading business has only five sources of
revenue.

\begin{method}[Classifying a trading revenue]\label{met:m1:what-a-trading-firm-does:sources}
Ask of any dollar of trading revenue which of these it is:
\begin{enumerate}
\item \textbf{Spread:} paid by someone who wanted to trade \emph{now}.
\item \textbf{Fee or commission:} paid for a service (managing, executing,
  clearing, lending, publishing data) whatever the market does.
\item \textbf{Risk premium:} paid, on average, for holding a risk others
  want to shed --- equity risk, illiquidity, the short side of an insurance.
\item \textbf{Carry and financing:} the difference between what a position
  yields and what it costs to fund.
\item \textbf{\omterm{def:m1:what-a-trading-firm-does:alpha}{Alpha}:} profit from being right about a price when others were
  wrong.
\end{enumerate}
A \omterm{def:m1:what-a-trading-firm-does:market-maker}{market maker} lives on 1; an \omterm{def:m1:what-a-trading-firm-does:asset-manager}{asset manager} on 2; an index investor collects
3; a bank's financing desk lives on 4; \omterm{def:m1:what-a-trading-firm-does:hedge-fund}{hedge funds} sell 5 and often deliver a
mixture of 3 and 5.
\end{method}

\begin{definition}[Alpha]\label{def:m1:what-a-trading-firm-does:alpha}
The \emph{alpha}\index{alpha} of a strategy is the part of its expected
return that is not compensation for bearing risks available cheaply
elsewhere: in a regression of its returns on those of the tradable risk
factors, the intercept.
\end{definition}

Source 1 looks like free money: buy at the bid, sell at the ask, repeat. It
is not, and the reason is the most important idea of the chapter.

\begin{definition}[Adverse selection]\label{def:m1:what-a-trading-firm-does:adverse-selection}
A \omterm{def:m1:what-a-trading-firm-does:market-maker}{market maker} suffers \emph{adverse selection}\index{adverse selection}
when the counterparties who choose to trade with its quotes are, more often
than chance, those who know the price is about to move against it. Its cost
is measured as the average move of the mid against the \omterm{def:m1:what-a-trading-firm-does:market-maker}{market maker}'s new
position after a trade.
\end{definition}

\begin{proposition}[Net capture]\label{prop:m1:what-a-trading-firm-does:capture}
A \omterm{def:m1:what-a-trading-firm-does:market-maker}{market maker} quotes $m \pm s/2$. A fraction $p$ of the orders it fills are
informed: after each, the mid moves by $J$ against the \omterm{def:m1:what-a-trading-firm-does:market-maker}{market maker}. The
other orders carry no information. With a rebate $\rho$ and a clearing fee
$\kappa$ per share, the expected profit per share traded, the
\emph{net capture}\index{net capture}, is
\[
c \;=\; \frac{s}{2} \;-\; pJ \;+\; \rho \;-\; \kappa .
\]
\end{proposition}
\begin{proof}
A customer buy of one share at $m + s/2$ leaves the \omterm{def:m1:what-a-trading-firm-does:market-maker}{market maker} short one
share at that price. Marked at the mid it has gained $s/2$. With probability
$p$ the mid then rises by $J$ and the short position loses $J$; otherwise the
expected move is zero. Sells are symmetric. Rebate and fee are paid per
share, whatever happens next. Positions inherited from earlier trades have
zero expected profit, because the side of the next order is independent of
them.
\end{proof}

\begin{example}[One cent is not one cent]\label{ex:m1:what-a-trading-firm-does:capture}
Take $s/2 = 1.00$ cent, $p = 15\%$, $J = 3$ cents, $\rho = 0.20$ cent and
$\kappa = 0.02$ cent. Then $c = 1.00 - 0.45 + 0.20 - 0.02 = 0.73$ cent: almost
half of the spread earned is handed back to better-informed counterparties.
If $p$ rises to $35\%$, $c = 0.13$ cent; at $40\%$ the business loses money
on every share while still ``earning the spread'' on each.
\end{example}

\begin{proposition}[Break-even volume]\label{prop:m1:what-a-trading-firm-does:breakeven}
A firm with fixed costs $C$ per day and \omterm{prop:m1:what-a-trading-firm-does:capture}{net capture} $c > 0$ per share breaks
even at a daily volume $V^\star = C/c$, and its daily profit at volume $V$ is
$c\,(V - V^\star)$.
\end{proposition}
\begin{proof}
Daily profit is $cV - C = c\,(V - C/c)$.
\end{proof}

The proposition is trivial and it is the whole economics of electronic
market making: costs are fixed (people, machines, data, connectivity) and
revenue is a tiny number multiplied by a very large one. Above $V^\star$
nearly every extra share is profit; below it the firm bleeds. This is why the
industry concentrates, and why public \omterm{def:m1:what-a-trading-firm-does:market-maker}{market makers}' results swing with
market volumes.

\begin{omfigure}
\begin{tikzpicture}
\begin{axis}[width=11cm,height=5cm,
  xlabel={daily volume (million shares)},ylabel={annual profit (\$ million)},
  tick label style={font=\small},label style={font=\small},
  xmin=0,xmax=60,ymin=-30,ymax=50,grid=major,grid style={gray!25}]
\addplot[omDef,very thick] table[col sep=comma,x=volume_m,y=profit_m]{figdata/markets-1/01-what-a-trading-firm-does/break_even.csv};
\addplot[omThm,dashed] coordinates {(20,-30) (20,50)};
\node[omThm,font=\small,anchor=west] at (axis cs:21,-10) {$V^\star = 20$ million};
\end{axis}
\end{tikzpicture}

\omcaption{Annual profit against daily volume for a \omterm{prop:m1:what-a-trading-firm-does:capture}{net capture} of 0.45 cent a
share and fixed costs of \$90\,000 a day. Below $V^\star$ the firm loses up
to its whole cost base; above it, profit grows without new cost. Data:
computed by the chapter's script.}
\end{omfigure}

\begin{dated}{2026-09}{Two listed market makers, in their own numbers}\label{dat:m1:what-a-trading-firm-does:listed}
Most \omterm{def:m1:what-a-trading-firm-does:market-maker}{market makers} are private, but two large ones publish results.
Virtu Financial reported total revenues of \$3\,632 million and adjusted net
trading income of \$2\,145 million for 2025 (\$1\,598 million for 2024).
Flow Traders, a specialist in exchange-traded products, reported net trading
income of \euro 486 million for 2025 with 635 employees at year end; for 2024
it reported \euro 468 million, 609 employees, and \euro 1\,545 billion of
exchange-traded products traded. Exercise~\ref{exo:m1:what-a-trading-firm-does:6}
turns these into a capture in basis points.
\end{dated}

\section{Principal and agent}

\begin{definition}[Principal and agent]\label{def:m1:what-a-trading-firm-does:principal-agent}
A firm trades as \emph{principal}\index{principal} when it is itself the
buyer or the seller: the position, and its risk, land on its own balance
sheet. It trades as \emph{agent}\index{agent} when it arranges a trade
between its client and a third party: it never owns the position and is paid
a commission.
\end{definition}

\begin{example}[The same sale, twice]\label{ex:m1:what-a-trading-firm-does:twice}
A client wants to sell 500\,000 shares trading around \$50. As
\emph{\omterm{def:m1:what-a-trading-firm-does:principal-agent}{agent}}, the broker works the order in the market over the day for a
commission of 1 cent a share: it earns \$5\,000 whatever price the client
obtains, and the client bears the risk that the price falls during the day.
As \emph{\omterm{def:m1:what-a-trading-firm-does:principal-agent}{principal}}, the dealer buys the whole block now at \$49.80: the
client is done, and the dealer owns \$24.9 million of stock that it must
sell. If it sells at an average of \$49.92 it makes \$60\,000; at \$49.70 it
loses \$50\,000. The 20-cent discount is the price of transferring the risk.
\end{example}

\begin{remark}[Why the distinction is policed]\label{rem:m1:what-a-trading-firm-does:conflict}
An \omterm{def:m1:what-a-trading-firm-does:principal-agent}{agent} owes its client the best result it can obtain; a \omterm{def:m1:what-a-trading-firm-does:principal-agent}{principal} is the
client's counterparty and profits from a worse price for the client. A firm
that does both --- most banks and most retail brokers' execution partners ---
must tell the client in which capacity it acted on each trade, and regulators
on both sides of the Atlantic examine how it chose. The rules are the subject
of \cref{ch:m1:retail-flow-and-wholesaling,ch:m1:european-equity-market-structure}.
\end{remark}

\section{A map of one trade}

\begin{omfigure}
\begin{tikzpicture}[font=\small,
  box/.style={draw=omDef,rounded corners=2pt,minimum width=2.5cm,minimum height=0.9cm,align=center,fill=omDef!4},
  svc/.style={draw=omExo,rounded corners=2pt,minimum width=2.5cm,minimum height=0.9cm,align=center,fill=omExo!6},
  flow/.style={-{Stealth[length=2mm]},thick,omDef},
  pay/.style={-{Stealth[length=2mm]},dashed,omThm}]
\node[box] (inv) at (0,0) {Investor};
\node[box] (brk) at (4.3,0) {Broker\\(agent)};
\node[box] (exc) at (8.6,0) {Exchange};
\node[box] (mm) at (12.9,0) {Market maker\\(principal)};
\node[svc] (ccp) at (8.6,-2.6) {Clearing house};
\node[svc] (dat) at (8.6,2.4) {Data vendors};
\draw[flow] (inv) -- node[above]{buy} (brk);
\draw[flow] (brk) -- node[above]{order} (exc);
\draw[flow] (mm) -- node[above]{ask} (exc);
\draw[flow] (exc) -- node[right]{trade} (ccp);
\draw[flow] (exc) -- node[right]{prices} (dat);
\draw[pay] (inv.south) to[bend right=25] node[below,pos=0.5]{commission} (brk.south);
\draw[pay] (brk.south east) to[bend right=20] node[below left,pos=0.75]{taker fee} (exc.south west);
\draw[pay] (exc.south east) to[bend right=20] node[below right,pos=0.3]{rebate} (mm.south west);
\draw[pay] (mm.south) to[bend left=25] node[below right,pos=0.5]{clearing fee} (ccp.east);
\end{tikzpicture}

\omcaption{One retail buy order. Solid arrows carry orders and information;
dashed arrows carry money other than the price of the shares. The half-spread
does not appear as an arrow: it is inside the price the investor pays.}
\end{omfigure}

Follow the dentist's hundred shares. Her broker, acting as \omterm{def:m1:what-a-trading-firm-does:principal-agent}{agent}, routes the
order to an exchange and charges a commission (often zero to her face, and
recovered elsewhere, as \cref{ch:m1:retail-flow-and-wholesaling} explains).
The exchange matches it against the \omterm{def:m1:what-a-trading-firm-does:market-maker}{market maker}'s resting ask; under the
common \emph{maker--taker} schedule it charges the side that took liquidity
and pays part of that to the side that provided it
(\cref{ch:m1:us-equity-market-structure}). The clearing house steps between
the two sides and guarantees settlement, for a fee
(\cref{ch:m1:clearing-and-settlement}). The exchange sells the record of the
trade to data vendors and directly to firms
(\cref{ch:m1:exchanges-brokers-venues}). And the \omterm{def:m1:what-a-trading-firm-does:market-maker}{market maker}, now short a
hundred shares it sold at \$50.01, waits for a seller to arrive at \$49.99 ---
and hopes the dentist knew nothing.

\section{Tutorial: who earned what today?}\label{tut:m1:what-a-trading-firm-does:whoearns}

\begin{tutorial}
\textbf{Goal.} Simulate one day of trading in one stock and split every
dollar that changed hands between the \omterm{def:m1:what-a-trading-firm-does:market-maker}{market maker}, the exchange, the broker
and the clearing house. \textbf{End state:} the bar chart that follows this tutorial,
and a check of
\cref{prop:m1:what-a-trading-firm-does:capture} against simulation.
\begin{tutsteps}
\item \textbf{Read the model.} Twenty thousand customer orders of 100 shares
hit a quote of $m \pm 1$ cent; 15\% of them are informed and move the mid by
3 cents. Fees follow a maker--taker schedule.
\omcode{code/markets-1/01-what-a-trading-firm-does/python/whoearns.py}{30}{41}{The trading loop: the market maker takes the other side of every order, and the informed move the price afterwards.}\label{lst:m1:what-a-trading-firm-does:loop}
\item \textbf{Split the \omterm{def:m1:what-a-trading-firm-does:market-maker}{market maker}'s result.} The position is marked at
the closing mid. Whatever is not spread earned is the profit or loss of
holding inventory while prices moved: mostly \omterm{def:m1:what-a-trading-firm-does:adverse-selection}{adverse selection}.
\omcode{code/markets-1/01-what-a-trading-firm-does/python/whoearns.py}{42}{55}{Who receives what. The customers' cost equals the spread plus the commission; the professionals then share it out.}\label{lst:m1:what-a-trading-firm-does:split}
\item \textbf{Run it.} From a Python prompt in the chapter's
\texttt{python/} directory, import \texttt{whoearns} and print
\texttt{simulate\_day(1)}. You should see a spread earned of \$20\,000, a position loss near \$5\,600
and a market-maker net near \$18\,000.
\item \textbf{Compare with theory.} The proposition predicts a position loss
of $pJ = 0.45$ cent a share, \$9\,000 on two million shares. One day is one
draw: run a thousand seeds and the mean converges to it, with a standard
deviation of about \$17\,000 a day (the histogram below).
\end{tutsteps}
\textbf{What to change next.} Set \texttt{informed=0.35}: the firm still
earns \$20\,000 of spread a day and now barely breaks even. Then make the
\omterm{def:m1:what-a-trading-firm-does:market-maker}{market maker} widen its quote after each informed trade, and watch volume ---
which the model holds fixed --- become the thing you would need to model.
\end{tutorial}

\begin{omfigure}
\begin{tikzpicture}
\begin{axis}[width=11.5cm,height=5.4cm,xbar,bar width=7pt,
  xlabel={\$ thousand, one simulated day},
  ytick=data,yticklabels from table={figdata/markets-1/01-what-a-trading-firm-does/whoearns.csv}{label},
  yticklabel style={font=\small},tick label style={font=\small},label style={font=\small},
  y dir=reverse,xmin=-10,xmax=25,enlarge y limits=0.1,grid=major,grid style={gray!25},
  table/col sep=comma]
\addplot[fill=omDef!60,draw=omDef] table[x=usd_thousands,y=k]{figdata/markets-1/01-what-a-trading-firm-does/whoearns.csv};
\end{axis}
\end{tikzpicture}

\omcaption{Who earned what on one simulated day of two million shares. The
\omterm{def:m1:what-a-trading-firm-does:market-maker}{market maker}'s net is the first three bars less its clearing fees. Data:
the tutorial's simulator, seed 1.}
\end{omfigure}

\begin{omfigure}
\begin{tikzpicture}
\begin{axis}[width=11cm,height=4.8cm,ybar,bar width=5.2pt,
  xlabel={market maker's daily net (\$ thousand)},ylabel={days out of 1\,000},
  tick label style={font=\small},label style={font=\small},
  xmin=-62,xmax=82,ymin=0,xtick={-60,-40,-20,0,20,40,60,80},table/col sep=comma]
\addplot[fill=omDef!40,draw=omDef] table[x=centre,y=count]{figdata/markets-1/01-what-a-trading-firm-does/daily_net_hist.csv};
\end{axis}
\end{tikzpicture}

\omcaption{A thousand simulated days. The mean is the \omterm{prop:m1:what-a-trading-firm-does:capture}{net capture} times the
volume; the spread around it is inventory risk. Six days in seven are
profitable. Data: the tutorial's simulator, seeds 0 to 999.}
\end{omfigure}

\section{Build: the firm's ledger}\label{bld:m1:what-a-trading-firm-does:ledger}

\begin{build}
\bfield{Purpose} The miniature trading firm assembled across this series
needs one place where every dollar earned or spent is recorded and
classified. Every later component --- the P\&L keeper of
\cref{ch:m1:pnl-and-positions}, the fee engines, the financing calculator ---
posts to it.
\bfield{Interface} An immutable \texttt{Entry(ts, account, category,
currency, amount, ref)} and a \texttt{Ledger} with \texttt{post(entry)},
\texttt{balance(account\_prefix, currency)} and
\texttt{by\_category(account\_prefix, currency)}. Accounts are dotted paths
(\texttt{desk.mm.equities}); a prefix selects a subtree. Categories are
\texttt{spread}, \texttt{position}, \texttt{fee}, \texttt{rebate},
\texttt{commission}, \texttt{financing}, \texttt{other}.
\bfield{Rules} Amounts are signed \emph{integers} in ten-thousandths of the
currency unit, so that a rebate of 0.2 cent is exactly 20 units: never
floating point for money that must add up. Entries are posted in time order;
unknown categories, malformed currencies and non-integer amounts are
rejected. Currencies are never added together.
\bfield{Acceptance tests} \texttt{code/firm/ledger/tests/test\_ledger.py}:
ten thousand sub-cent rebates sum exactly; prefixes and categories aggregate;
bad entries raise.
\bfield{Stretch} Post the tutorial's simulated day trade by trade and
reproduce the tutorial's bar chart from the ledger alone.
\end{build}

\begin{omsources}
\item Virtu Financial, \emph{Fourth Quarter 2025 Results}, press release,
29 January 2026; \emph{Fourth Quarter 2024 Results}, 29 January 2025.
\item Flow Traders, \emph{4Q and FY 2025 Results}, 12 February 2026;
\emph{4Q and FY 2024 Results}, 13 February 2025.
\item Boston Consulting Group, \emph{Global Asset Management Report 2026};
\emph{Global Asset Management Report 2025}.
\item ``Alfred Winslow Jones'', \emph{Institutional Investor}, Hall of Fame
profile.
\item L.~Harris, \emph{Trading and Exchanges}, Oxford University Press, 2003,
chapters 3, 13 and 14 --- the classic taxonomy of traders and of dealers'
revenues.
\end{omsources}

\section{Exercises}

\begin{exercise}[$\star$]\label{exo:m1:what-a-trading-firm-does:1}
A stock is quoted $24.98$ bid, $25.02$ ask. Give the spread, the mid, the
spread in basis points of the mid, and the cost relative to the mid of a
market order for 300 shares.
\end{exercise}

\begin{exercise}[$\star$]\label{exo:m1:what-a-trading-firm-does:2}
A \omterm{def:m1:what-a-trading-firm-does:market-maker}{market maker} trades 2 million shares a day with a \omterm{prop:m1:what-a-trading-firm-does:capture}{net capture} of 0.1 cent
a share. What does it earn in a day, and in a year of 252 trading days?
\end{exercise}

\begin{exercise}[$\star$]\label{exo:m1:what-a-trading-firm-does:3}
For each activity say whether the firm acts as \omterm{def:m1:what-a-trading-firm-does:principal-agent}{principal} or as \omterm{def:m1:what-a-trading-firm-does:principal-agent}{agent}, and
which of the five revenue sources of
\cref{met:m1:what-a-trading-firm-does:sources} pays it: (a) a broker routes
a client's order to an exchange for a commission; (b) a dealer buys a block
from a client at a discount; (c) a manager runs a pension fund's equity
portfolio for 0.15\% a year; (d) a firm quotes both sides of an
exchange-traded fund all day; (e) a fund buys a stock because its model
predicts a rise.
\end{exercise}

\begin{exercise}[$\star\star$]\label{exo:m1:what-a-trading-firm-does:4}
A \omterm{def:m1:what-a-trading-firm-does:hedge-fund}{hedge fund} starts the year with \$500 million and earns 12\% before fees.
It charges a 2\% management fee on opening assets and 20\% of the profit
remaining after that fee. Compute both fees, the investors' net return, and
the share of the gross profit that went to the manager.
\end{exercise}

\begin{exercise}[$\star\star$]\label{exo:m1:what-a-trading-firm-does:5}
A \omterm{def:m1:what-a-trading-firm-does:market-maker}{market maker} earns a half-spread of 1 cent; 14\% of its fills are
informed and are followed by a 5-cent adverse move; it receives a rebate of
0.20 cent and pays 0.05 cent in clearing per share. Compute its \omterm{prop:m1:what-a-trading-firm-does:capture}{net capture}.
Its fixed costs are \$90\,000 a day: find its break-even volume.
\end{exercise}

\begin{exercise}[$\star\star$]\label{exo:m1:what-a-trading-firm-does:6}
Use the 2024 figures of
\cref{dat:m1:what-a-trading-firm-does:listed}. (a) Express Flow Traders' net
trading income as basis points of the value of exchange-traded products it
traded. Why is this an upper bound on its capture in those products?
(b) Compute net trading income per employee. (c) On a \euro 50 product, how
many cents is the capture of (a)?
\end{exercise}

\begin{exercise}[$\star\star\star$]\label{exo:m1:what-a-trading-firm-does:7}
\emph{Coding.} With the tutorial's simulator, run \texttt{simulate\_day(7)}.
(a) Report the \omterm{def:m1:what-a-trading-firm-does:market-maker}{market maker}'s spread earned, position P\&L and net. (b) The
position P\&L is far from its expectation of $-\$9\,000$ on many days.
Without running anything, explain which term of the model creates this
dispersion, and how it would scale if each order were for 400 shares and
there were 5\,000 of them.
\end{exercise}

\begin{exercise}[$\star\star\star$]\label{exo:m1:what-a-trading-firm-does:8}
\emph{Find the flaw.} A pitch deck reads: ``Our strategy earns the \omterm{def:m1:what-a-trading-firm-does:spread}{bid--ask
spread}. The average spread in our universe is 6~basis points, we trade our
capital ten times a day, so we earn 60~basis points a day before costs. We
execute with market orders to guarantee our fills.'' Find two independent
errors, and say what the daily figure is likely to be.
\end{exercise}

\section{Problem: A Week at Quad Street Markets}

\begin{problem}[{Weekend problem --- a week at a small market maker}]\label{pb:m1:what-a-trading-firm-does:1}
Quad Street Markets is a fourteen-person proprietary firm making markets in
forty large stocks. Its owner asks you one question: ``how much do we have to
trade to survive?''

\medskip\noindent\textbf{Part I --- What a share earns.} The average stock
trades at \$40 with a spread of 2 cents. Quad Street trades 6 million shares
a day, always passively.
\begin{enumerate}
\item Give the half-spread in cents and in basis points of the price.
\item 12\% of Quad Street's fills are informed and are followed by a 5-cent
adverse move. What does \omterm{def:m1:what-a-trading-firm-does:adverse-selection}{adverse selection} cost per share?
\item The exchanges pay a rebate of 0.20 cent; clearing costs 0.02 cent.
Compute the \omterm{prop:m1:what-a-trading-firm-does:capture}{net capture}.
\item Compute the expected daily trading revenue.
\item And the expected annual revenue, for 252 days.
\end{enumerate}

\medskip\noindent\textbf{Part II --- What the firm costs.}
\begin{enumerate}[resume]
\item Fourteen people cost \$350\,000 each per year all-in; technology,
colocation and data cost \$2.1 million; everything else \$0.56 million. Give
the annual cost and the cost per trading day.
\item Give the expected daily and annual profit.
\item Give the cost-to-income ratio.
\end{enumerate}

\medskip\noindent\textbf{Part III --- Breaking even.}
\begin{enumerate}[resume]
\item Find the break-even daily volume $V^\star$.
\item A new aggressive fund becomes active in Quad Street's stocks and the
informed share of its fills rises to 16\%. Find the new \omterm{prop:m1:what-a-trading-firm-does:capture}{net capture} and the
new $V^\star$.
\item At 6 million shares a day, what is the daily result now?
\item What half-spread would restore the original \omterm{prop:m1:what-a-trading-firm-does:capture}{net capture}? The minimum
price increment is 1 cent: what can Quad Street actually do?
\item Independently of question 10, the exchanges cut the rebate to 0.15
cent. Find $V^\star$.
\end{enumerate}

\medskip\noindent\textbf{Part IV --- Risk.} Return to the original figures.
The daily result has a standard deviation of \$22\,000 and days are
independent and roughly normal.
\begin{enumerate}[resume]
\item What is the probability of a losing day?
\item Compute the annualised Sharpe ratio (mean over standard deviation of
the daily result, times $\sqrt{252}$).
\item What is the probability of a losing year?
\item The owners keep \$5 million of capital in the firm. What is the
expected annual return on it?
\item Volumes double in a volatile quarter and costs do not move. By what
factor does the daily profit grow?
\item Explain in two sentences why a \omterm{def:m1:what-a-trading-firm-does:market-maker}{market maker}'s owners like volatile
markets and its risk manager does not.
\item State the answer to the owner's question in one sentence, naming the
three numbers he should watch every day.
\end{enumerate}
\end{problem}

\section{Interview questions}

\begin{interviewq}[$\star$ \iqroles{trader, researcher, developer}]\label{iq:m1:what-a-trading-firm-does:1}
How does a \omterm{def:m1:what-a-trading-firm-does:market-maker}{market maker} make money, and how does it lose money?
\end{interviewq}

\begin{interviewq}[$\star$ \iqroles{trader, researcher, developer, mle}]\label{iq:m1:what-a-trading-firm-does:2}
What is the difference between a \omterm{def:m1:what-a-trading-firm-does:prop}{proprietary trading firm} and a \omterm{def:m1:what-a-trading-firm-does:hedge-fund}{hedge fund}?
Why might the same strategy be run differently in each?
\end{interviewq}

\begin{interviewq}[$\star\star$ \iqroles{trader, researcher}]\label{iq:m1:what-a-trading-firm-does:3}
You are quoting $49.99$ / $50.01$ and someone sells you 100 shares at
$49.99$. A minute later the market is $49.96$ / $49.98$. How much have you
made or lost, and how would you split that number into two parts?
\end{interviewq}

\begin{interviewq}[$\star\star$ \iqroles{researcher, developer}]\label{iq:m1:what-a-trading-firm-does:4}
A firm nets one basis point on the \$2 billion it trades each day. Estimate
its annual trading revenue in your head. Is one basis point a plausible
number?
\end{interviewq}

\begin{interviewq}[$\star\star$ \iqroles{bank, trader}]\label{iq:m1:what-a-trading-firm-does:5}
A client calls to sell a block worth ten days of average volume. Describe
two ways your desk can help, how the desk is paid in each, and who bears the
price risk.
\end{interviewq}

\begin{interviewq}[$\star\star\star$ \iqroles{trader, researcher}]\label{iq:m1:what-a-trading-firm-does:6}
Last month the spreads you quoted were twice as wide as the month before,
your volume was the same, and you made less money. Give two explanations and
the measurement that would tell them apart.
\end{interviewq}
